\begin{claim}{B.26}{Proved}{Complete executed spherical catalogues through six punctures}
For $\nu=3,4,5,6$, the extended mapping class group of the sphere with
$\nu$ unlabelled punctures has respectively
\[
                         10,\quad30,\quad20,\quad57
\]
conjugacy classes of finite subgroups, including the trivial subgroup.
The exact finite parameterization below gives every class once.  The
executed files \texttt{type\_0\_\nu\_orbifold\_catalogue.csv} and
\texttt{type\_0\_\nu\_orbifold\_sections.json}, for each of these four
values of $\nu$, give all partial-section classes, all their inclusion
relations, $\Phi$, its group action and stabilizers, the centralizers,
the multiplier on each section element, and the puncture stabilizers.
Thus these files settle the whole computation requested in B(f) for
the hyperbolic spherical types with $2g+\nu\le6$.
\end{claim}
\begin{proof}
Extend an action over the punctures.  By B.1 and the sphere theorem of
\cite[Theorem 1.2]{bKolev2006}, a finite topological action on $S^2$ is
conjugate to an orthogonal action.  Write $Q=S^2/G$ with its local
orbifold groups.  A point has cyclic positive stabilizer, or a dihedral
stabilizer with mirrors; the mirror-only stabilizer has order two.
The covering formula is
\[
 \chi_{\rm orb}(Q)=\chi(|Q|)
  -\sum_{j=1}^r(1-1/m_j)
  -\frac12\sum_{j=1}^k(1-1/n_j)=\frac2N,
 \qquad N=|G|.                                      \tag{B.26.1}
\]
Here the $m_j$ are interior cone orders and the $n_j$ are corner orders.
One proves this formula directly by a triangulation whose vertices
include all cone and corner points: count each cell in the covering
with weight equal to the reciprocal of its stabilizer order.  Mirrors
are boundary strata; they introduce no additional term after the
boundary cells and their endpoints are counted together.

Since (B.26.1) is positive, the underlying compact surface is $S^2$,
$\mathbb{RP}^2$, or a disc.  The sphere has at most three cones.  A
single cone is impossible, since its local generator would be trivial
in the orbifold presentation.  Two cones have equal orders, since the
two local generators are inverses and both local cyclic groups inject.
For three cones the inequality $1/m_1+1/m_2+1/m_3>1$ gives precisely
$(2,2,n),(2,3,3),(2,3,4),(2,3,5)$.
The projective plane has at most one cone.  For a disc, positivity gives
$2r+k\le3$.  A single corner without a cone is impossible; two corners
without a cone have equal orders, by the same local-injectivity argument.
Three corners give the same four triples.  One cone and one corner
give $2/m+1/n>1$, hence $(m,n)=(2,n)$ or $(3,2)$.
Consequently the following table is exhaustive.  In its notation
$S(m_1,\ldots,m_r)$ is a sphere with cones, $P(m)$ is a projective
plane with one cone, $D(m_1,\ldots,m_r;n_1,\ldots,n_k)$ is a mirror
disc with the specified interior cones and cyclically ordered corners,
and an empty list means no such singularities.
\[
\begin{array}{c|c@{\qquad}c|c}
Q&N&Q&N\\\hline
S()&1&P()&2\\
S(n,n)&n&P(n)&2n\\
S(2,2,n)&2n&D(;)&2\\
S(2,3,3)&12&D(;n,n)&2n\\
S(2,3,4)&24&D(n;)&2n\\
S(2,3,5)&60&D(;2,2,n)&4n\\
&&D(2;n)&4n\\
&&D(;2,3,3)&24\\
&&D(;2,3,4)&48\\
&&D(;2,3,5)&120\\
&&D(3;2)&24
\end{array}                                                   \tag{B.26.2}
\]
Every integer $n$ in this table is at least two.  For $3\le\nu\le6$
only $2\le n\le6$ can occur.  Indeed, in every family depending on
$n$, the union of all point orbits of size strictly smaller than $n$
has at most two points; an invariant set of at least three punctures
therefore contains an orbit of size at least $n$.  This assertion also
follows immediately by dividing $N$ by each local group order in
(B.26.2).  In the two icosahedral rows every point orbit has at least
twelve points.  Those rows cannot occur in the present range.  These
are proved finite bounds, rather than bounds inferred from the output.

We now give the exact decoration rule used in the program.  A cone
orbit may be removed or retained; put $c_j\in\{0,1\}$ accordingly.
A corner orbit similarly has $d_j\in\{0,1\}$.  For each open mirror
edge let $e_i\ge0$ be the number of removed orbits on that edge; a
cornerless mirror circle is one edge for this purpose.  Finally let
$a\ge0$ be the number of regular interior orbits removed.  The single
equation is
\[
 \sum_jc_j\frac N{m_j}
 +\sum_jd_j\frac N{2n_j}
 +\sum_i e_i\frac N2+aN=\nu.                       \tag{B.26.3}
\]
On a sphere, permutations of cones of equal order identify decorations.
On a mirror disc, identify them under precisely the rotations and
reflections of the cyclic corner list that preserve its order labels.
The corresponding edge counts move with their edges.  For two corners
this rule includes the independent exchanges of the two corners and
the two edges.  For one or zero corners it introduces no further
identification of the single edge count.  There is at most one interior
cone in every remaining disc or projective-plane row.

For clarity, this rule also proves completeness of the conjugacy
classification.  A homeomorphism of an interval moves any two
unlabelled finite subsets of the same size into one another, and the
same assertion holds in the connected regular interior of each quotient
surface.  Such moves extend to the quotient surface by disjoint small
discs and strips.  The only remaining data are therefore the selected
singular strata and the cyclic mirror-edge counts in (B.26.3), modulo
the explicitly stated symmetries.  Conversely every quotient
homeomorphism preserves those strata, their orders, and their cyclic
incidence.  Since $S^2$ is the universal orbifold cover, a homeomorphism
between the decorated quotient orbifolds lifts to a conjugacy between
their covering actions; this follows by lifting paths and their local
cone or mirror charts.  A conjugacy of actions descends to such a
quotient homeomorphism.  Thus no extra equivalence of decorations is
possible.  Finally C.5 and C.6 identify these topological action classes
with conjugacy classes of finite subgroups in the extended mapping
class group of the punctured sphere.  The induced action there is
faithful: a nonidentity positive periodic map isotopic to the identity
would fix every puncture individually, whereas a nonidentity sphere
rotation has exactly two fixed points; a negative map cannot be
isotopic to the identity.  Here $\nu\ge3$ is essential.

Here are certificates for the finite group calculations.  The source
\texttt{catalogue\_spheres.py} constructs the cyclic and dihedral
families as regular permutation groups, the tetrahedral and octahedral
positive groups as $A_4,S_4$, their needed extensions by an explicitly
given central involution, and the full octahedral group as signed
permutations of three axes.  All generators, products, orders and signs
are recorded as finite permutations.  The program verifies each
orbifold relation, injection of every local cyclic or dihedral group,
and $N\chi_{\rm orb}(Q)=2$ using rational arithmetic.  These checks
prove that the permutation group is the entire orbifold group, as
follows.  The kernel of the displayed quotient of the orbifold
presentation produces a covering with no cone, corner or mirror
singularity because all local groups inject.  Its orientation
character is trivial, and its Euler characteristic is two.  It is
therefore a closed orientable sphere; its fundamental group is trivial.
This proves injectivity of the quotient presentation as well as its
already checked surjectivity.  In particular these certificates do
not merely exhibit a possibly proper finite quotient of an unexamined
presentation.

We describe how the section data are recovered without numerical
sampling.  Cone and corner lifts are cosets of their cyclic and
dihedral stabilizers.  In a disc presentation choose reflections
$r_0,\ldots,r_k$ with
\[
 (r_jr_{j+1})^{n_j}=1,\qquad
 r_k=e r_0e^{-1},\qquad x_1\cdots x_r e=1.            \tag{B.26.4}
\]
Indices of the corner orders are shifted here to start at zero.
The code stores $r_0,\ldots,r_{k-1}$ and derives the final reflection.
An edge lift is a coset of $\langle r_i\rangle$ and meets the adjacent
corner cosets; the closing incidence at edge zero is multiplied by
$e^{-1}$.  With no corners it instead joins the corresponding two
cosets of $\langle r_0\rangle$ by that same rule.  Insert the prescribed
removed vertices along each edge, retaining two auxiliary vertices in
each complementary segment.  These auxiliary vertices prevent a
one-edge loop from collapsing when represented as an ordinary finite
adjacency graph.  This is exactly the lifted mirror locus with a
finite subdivision.

For any reflection $r\in G$, keep the vertices whose stabilizers
contain $r$.  Before punctures are deleted the program verifies degree
two at every such vertex.  Its components are consequently the full
fixed circles of $r$.  Deleting the selected puncture vertices gives
all the fixed lines and ovals, with their exact incidence.  For a
subgroup whose positive part is nontrivial, the fixed points are
precisely the retained cone and corner lifts whose stabilizers contain
that subgroup.  A negative involution without a mirror has an empty
fixed set.  Thus B.3 and B.4 supply every section class, with centralizer
$1$ for a point or a line and $\Z$ for an oval.  A section over $H$
extends a section over $K<H$ exactly when its fixed component is
contained in the latter fixed component.  Hence the same finite
incidence calculation gives the entire partial-section poset and all
its maximal elements, including the reflection cases with several
maximal extensions.  Left multiplication on cosets gives the group
action and every stabilizer without a choice of orbit representatives.

The multiplier is exact as well.  Orient a local lift of each cone,
and at a corner use the positive generator $r_jr_{j+1}$ determined by
the oriented boundary order.  Write this generator as $u$, of order
$m$.  At the point represented by a coset $gD$, an element $h$ of its
positive stabilizer satisfies $g^{-1}hg=u^b$, and
\[
                   \mu(h)=\operatorname{sgn}(g)\frac bm
                                \pmod{\Z}.           \tag{B.26.5}
\]
Replacing $g$ by another coset representative leaves this value
unchanged: a positive local element commutes with $u$, while a
negative local element both reverses the sign and inverts $u$.
For a reflection the value is $1/2$, and for the identity it is zero.
Formula (B.26.5) is the local instance of B.12, and applies to all
sections in the files.  The code checks its order and every power
identity, as well as uniqueness of the maximal extension for every
section with nontrivial positive part.  All nontrivial subgroups of
each finite permutation group are enumerated by saturation under
adjoining each group element; this stops only when the subgroup set
is closed under every possible enlargement.

The following compact census is part of the executed output; each
entry is the number of conjugacy classes of that group order.  A dash
denotes zero.  Complete individual rows, rather than just this census,
are in the CSV files specified in the statement.
\[
\begin{array}{c|rrrrrrrrrrrrr|r}
\nu\backslash N&1&2&3&4&5&6&8&10&12&16&20&24&48&\text{total}\\\hline
3&1&3&1&1&-&3&-&-&1&-&-&-&-&10\\
4&1&6&1&11&-&1&7&-&1&1&-&1&-&30\\
5&1&4&1&4&1&3&1&3&1&-&1&-&-&20\\
6&1&7&1&14&1&9&8&1&9&1&-&4&1&57
\end{array}                                                    \tag{B.26.6}
\]
As an independent check, the program compares the multisets of all
six columns consisting of group order, positive order, full-section
count, nontrivial partial-section count, maximal-section count, and
group-orbit count, for $\nu=3,4$, against the previous permutation and
amalgam calculations in B.23 and B.25.  Both multisets agree, with
multiplicities.  The check and its inputs are attached; the completeness
proof above does not depend on that agreement.
\end{proof}
